% Compile exam_01_reference.tex with LuaLaTeX or XeLaTeX.
% All content is supplied reference material; there are no student blanks.
% Sources: supplied FNRP Chapters 1--3, lesson handouts, and HW01--HW13.
\providecommand{\dd}{\,\mathrm d}
\newcommand{\is}[2]{\ensuremath{{}^{#1}\mathrm{#2}}}
\newcommand{\EE}{\mathbb E}
\newcommand{\fv}{\symbfit\phi}
\newcommand{\M}{\mathbf M}
\newcommand{\F}{\mathbf F}
\newcommand{\sv}{\mathbf s}
\newcommand{\RefHead}[3]{%
  \par\addvspace{3.5pt}%
  \noindent\colorbox{KSUPurple}{\color{white}\sffamily\bfseries\strut\hspace{1pt}#1\hspace{1pt}}%
  \hspace{0.35em}{\sffamily\bfseries\color{KSUDeepPurple}#2}%
  \hfill{\sffamily\scriptsize\color{KSUGray}#3}\par\vspace{1.5pt}}
\newcommand{\RefLabel}[1]{\textbf{\color{KSUDeepPurple}#1}}
\AtBeginDocument{%
  \fontsize{8.6}{9.6}\selectfont
  \setlength{\abovedisplayskip}{2pt plus .3pt minus .3pt}
  \setlength{\belowdisplayskip}{2pt plus .3pt minus .3pt}
  \setlength{\abovedisplayshortskip}{1.5pt}
  \setlength{\belowdisplayshortskip}{1.5pt}
  \setlength{\jot}{1.5pt}
  \setlength{\tabcolsep}{3.5pt}
  \setlength{\aboverulesep}{1.2pt}
  \setlength{\belowrulesep}{1.2pt}
}
\tcbset{handout base/.append style={
  before skip=3pt,after skip=3pt,top=3pt,bottom=3pt,
  fonttitle=\sffamily\bfseries\fontsize{8.6}{10}\selectfont,
  before upper={\RaggedRight}}}
\HandoutSetup{NE 630}{1--14}{Exam 1 Reference}{Principles, equations, and course data}
\fancyhead[R]{\sffamily\bfseries\small Lessons 1--14}
\begin{document}
\thispagestyle{handoutfirst}
\vspace*{\dimexpr-\headheight-\headsep\relax}
\begin{tcolorbox}[enhanced,frame hidden,colback=white,left=6pt,right=0pt,
  top=1pt,bottom=1pt,borderline west={4pt}{0pt}{KSUPurple},
  before skip=0pt,after skip=4pt]
  {\sffamily\bfseries\fontsize{17}{19}\selectfont Exam 1 Reference}
  \hfill {\sffamily\bfseries\large Lessons 1--14}\par
  {\sffamily\small NE 630\quad Nuclear reactions, neutron interactions, and infinite-medium balance}
\end{tcolorbox}
\begin{fixedbox}{Constants and units}
\begin{tabularx}{\linewidth}{@{}XXX@{}}
$N_A=6.022\times10^{23}\unit{mol^{-1}}$ &
$1\unit{b}=10^{-24}\unit{cm^2}$ &
$1\unit{eV}=1.60218\times10^{-19}\unit{J}$\\
$m_n=1.0086649\unit{u}=1.67493\times10^{-27}\unit{kg}$ &
$m_nc^2=939.57\unit{MeV}$ &
$c=2.9979\times10^8\unit{m/s}$\\
$k_B=8.6173\times10^{-5}\unit{eV/K}$ &
$1\unit{u}\,c^2=931.5\unit{MeV}$ &
$1\unit{Ci}=3.7\times10^{10}$ decays/s
\end{tabularx}
\end{fixedbox}
\begin{handoutcolumns}
% ---------- PAGE 1, LEFT ----------
\RefHead{1}{Reactions, energy, and power}{L1--2}
For ${}^{A}_{Z}X$, $A=Z+N$. Conserve total $A$, charge, energy,
and momentum; $Q>0$ releases kinetic/excitation energy.
\[
 Q=\Big(\sum m_{\rm in}-\sum m_{\rm out}\Big)c^2,
 \qquad B=\big[Z m({}^{1}{\rm H})+Nm_n-m_{
 X}^{\rm atom}\big]c^2.
\]
Greater $B/A$ means tighter binding. With \emph{atomic} masses,
$Q_{\beta^-}=(m_P-m_D)c^2$; do not subtract $m_e$ again.
A $\beta^-$ decay keeps $A$ fixed and increases $Z$ by one.
\[
 K_{\rm rel}=m_nc^2\left[(1-v^2/c^2)^{-1/2}-1\right],
 \qquad K_{\rm cl}=\tfrac12m_nv^2,
\]
\[
 \frac{v^2}{c^2}=1-\left(1+\frac{K_{\rm rel}}{m_nc^2}\right)^{-2},
 \quad \%\mathrm{err}=100\frac{K_{\rm cl}-K_{\rm rel}}{K_{\rm rel}}.
\]
Classically, $v=\sqrt{2E/m_n}$ (use consistent units).


\[
 N=\frac{mN_A}{M},\quad E_{\rm released}=N_{\rm reactions}Q,
 \quad \frac{\Delta m}{m}=\frac{E_{\rm released}}{mc^2},
\]
\[
 P_{\rm th}=E_f\dot N_f=E_fVR_f,
 \qquad P_e=\epsilon_{\rm th}P_{\rm th},\qquad E=P\Delta t.
\]
Use $E_f\simeq200\unit{MeV}=3.204\times10^{-11}\unit{J}$ per
fission unless specified. Prompt energy: fragments, neutrons, and
$\gamma$ rays. Later: $\beta^-$, $\gamma$, delayed neutrons, and
escaping antineutrinos. Fission products are usually neutron rich.
\textbf{Fissile}: low-energy fission (U-233, U-235, Pu-239).
\textbf{Fertile}: capture/decay produces fissile material
(Th-232 $\to$ U-233; U-238 $\to$ Pu-239).

\RefHead{2}{Multiplication and radioactive decay}{L2--3}
For fixed $k$ and generation time $\ell$,
\[
 n_i=n_0k^i,\quad
 \frac{n(t)}{n_0}=\frac{P(t)}{P_0}=k^{t/\ell}
 \simeq e^{(k-1)t/\ell}\quad(k\simeq1).
\]
$k<1$: subcritical; $k=1$: critical; $k>1$: supercritical.

\vspace{0.25cm}
\RefLabel{Decay heat after operation at $P_0$ for $t_0$.}
\[
 P_d(t)=0.0622P_0\big[t^{-0.2}-(t+t_0)^{-0.2}\big].
\]
Use $t,t_0$ in seconds; $t>0$ is time after shutdown.

\vspace{0.25cm}
\RefLabel{Inventory and activity.} $N$ is a number of nuclei; production
$R$ and activity $\mathcal A$ are nuclei/s and decays/s.
\[
 \lambda=\frac{\ln2}{t_{1/2}}=\frac1{\bar t},\qquad
 \mathcal A=\lambda N,\qquad m=\frac{\mathcal A M}{\lambda N_A}.
\]
\[
 \dot N=R-\lambda N\ \Longrightarrow\
 N(t)=N_0e^{-\lambda t}+\frac{R}{\lambda}(1-e^{-\lambda t})
 \quad(R\text{ constant}).
\]
For $R=0$, $N(nt_{1/2})=N_0/2^n$; for constant production,
$N_\infty=R/\lambda$ and $\mathcal A_\infty=R$. After production
stops, decay from the shutdown inventory.
For a one-way chain,
\[
 \dot N_i=R_i(t)+\lambda_{i-1}N_{i-1}-\lambda_iN_i,
\]
\[
 N_i(t)=N_i(0)e^{-\lambda_it}
 +\int_0^t e^{-\lambda_i(t-\tau)}
 [R_i(\tau)+\lambda_{i-1}N_{i-1}(\tau)]\dd\tau.
\]
For $A\to B\to C$, no production, $N_A(0)=N_0$, $N_B(0)=0$:
\[
 N_A=N_0e^{-\lambda_A t},\quad
 N_B=\frac{\lambda_AN_0}{\lambda_B-\lambda_A}
       (e^{-\lambda_A t}-e^{-\lambda_B t}).
\]
If $C$ is stable, $N_C=N_0-N_A-N_B$. Secular equilibrium:
$\lambda_A\ll\lambda_B$, after the daughter transient,
$\lambda_BN_B\simeq\lambda_AN_A$ (equal activities, not inventories).


\switchcolumn\RaggedRight
% ---------- PAGE 1, RIGHT ----------
\RefHead{3}{Composition, cross sections, and rates}{L4--5}
$M_i$: molar mass (g/mol); $n_i$: target density (cm$^{-3}$);
$x_i,w_i,f_{V,i}$: atom, mass, and volume fractions.
\[
 \overline M=\sum_i x_iM_i,\qquad
 n_i=x_i\frac{\rho N_A}{\overline M}
       =\frac{\rho w_iN_A}{M_i},\qquad
 x_i=\frac{w_i/M_i}{\sum_jw_j/M_j}.
\]
For a compound, $M_m=\sum_i a_iM_i$, $n_m=\rho N_A/M_m$,
$n_i=a_in_m$, with $a_i$ atoms per formula unit.


\[
 \Sigma_x=\sum_i n_i\sigma_x^i,\qquad
 \Sigma_{x,\rm mix}=\sum_i f_{V,i}\Sigma_{x,i}.
\]
Use $\sigma$ in cm$^2$ to obtain $\Sigma$ in cm$^{-1}$.
$\Sigma_t=\Sigma_s+\Sigma_a$; $\Sigma_a=\Sigma_\gamma+\Sigma_f$
when capture and fission are the only absorption channels.
\[
 R_x=\Sigma_x\phi\quad ,\qquad
 \overline\nu_{\rm mix}=\frac{\sum_i\nu_iR_f^i}{\sum_iR_f^i}.
\]
\RefLabel{Conditional probabilities at fixed energy.}
\[
 P(x\mid\mathrm{collision})=\frac{\Sigma_x}{\Sigma_t},\quad
 P(i\mid x)=\frac{n_i\sigma_x^i}{\Sigma_x},\quad
 P(f\mid a)=\frac{\Sigma_f}{\Sigma_a}.
\]
\RefLabel{Uncollided fluxes.}

Parallel beam through homogeneous slab:
\[
\phi_u(x)=\phi_u(0)e^{-\Sigma_t x}
\]
Isotropic point source $s_p$ (n/s) surrounded by multiple layers of thickness $\Delta_j$:
\[
 \phi_u(r)=\frac{s_p}{4\pi r^2}e^{-\sum_j\Sigma_{t,j} \Delta_j},
\]
for $\sum_j \Delta_j \leq r$.
Attenuation describes \emph{uncollided} flux, not total flux!

\RefHead{4}{Cross-section shapes and resonances}{L6}
Low-energy $1/v$: $\sigma_x(E)=\sigma_x(E_*)\sqrt{E_*/E}$.
Thresholds open channels; resonances reflect compound-nucleus levels.

\RefLabel{Single-level Breit--Wigner capture}
\[
 \Gamma=\Gamma_n+\Gamma_\gamma,\quad
 \lambda_r=\frac{4.55\times10^{-10}}{\sqrt{E_r/\mathrm{eV}}}\unit{cm},
 \quad \sigma_0=4\pi\lambda_r^2\frac{\Gamma_n}{\Gamma},
\]
\[
 \sigma_\gamma(E)=\sigma_0\frac{\Gamma_\gamma}{\Gamma}
 \sqrt{\frac{E_r}{E}}\,
 \frac{1}{1+4(E-E_r)^2/\Gamma^2}.
\]
The FWHM
is $\Gamma$. Doppler broadening lowers/widens peaks, approximately
preserving area.

\RefHead{5}{Elastic collisions and moderation}{L7--8}
\textbf{Stationary target, nonrelativistic, isotropic CM.}
$E\to E'$; $A=M/m_n\simeq$ mass number;
$\mu_c=\cos\theta_c$, $\mu_L=\cos\theta_L$.
\[
 \alpha=\left(\frac{A-1}{A+1}\right)^2,\quad
 \frac{E'}E=\frac{A^2+2A\mu_c+1}{(A+1)^2},\quad
 \mu_L=\frac{1+A\mu_c}{\sqrt{A^2+2A\mu_c+1}}.
\]
\[
 p(E\to E') = p(E'\mid E)=\frac{1}{(1-\alpha)E}\quad(\alpha E\le E'\le E);
 \qquad p=0\text{ otherwise}.
\]
\[
 \langle E'\rangle=\tfrac12(1+\alpha)E,\quad
 \langle E-E'\rangle=\tfrac12(1-\alpha)E,\quad
 \langle\mu_L\rangle=\frac{2}{3A}.
\]
\[
 \xi=\left\langle\ln\frac{E}{E'}\right\rangle
 =1+\frac{\alpha\ln\alpha}{1-\alpha},\quad
 n_{\rm coll}\simeq\frac{\ln(E_0/E_n)}{\xi}\quad(\xi\text{ constant}).
\]
For H-1, $\alpha=0$, $\xi=1$. Mixtures use \emph{collision} weights:
$p_{\rm mix}=\sum_i(\Sigma_s^i/\Sigma_s)p_i$ and
$\bar\xi=\sum_i\Sigma_s^i\xi_i/\Sigma_s$.
Slowing-down power $=\bar\xi\Sigma_s$; ratio $=\bar\xi\Sigma_s/\Sigma_a$.
\end{handoutcolumns}

\newpage
\begin{handoutcolumns}
% ---------- PAGE 2, LEFT ----------
\RefHead{6}{Densities, spectra, and expectations}{L7,9,11}
For $p\ge0$, $\int p\dd x=1$:
$P(a<X<b)=\int_a^bp\dd x$ and $\EE[f]=\int fp\dd x$.
A one-to-one transformation gives $p_Y(y)=p_X(x(y))|\dd x/\dd y|$.
Restrict every integral to the support.
\[
 \phi(E)=v(E)n(E),\quad \phi_{\rm tot}=\int\phi(E)\dd E,
 \quad R_x=\int\Sigma_x(E)\phi(E)\dd E.
\]
Units: $n(E)$ [cm$^{-3}$ eV$^{-1}$], $\phi(E)$ [cm$^{-2}$ s$^{-1}$ eV$^{-1}$],
$R_x$ [cm$^{-3}$ s$^{-1}$]; $v$ in cm/s.
\RefLabel{Lethargy.} $u=\ln(E_0/E)$, $\dd u=-\dd E/E$,
$\phi_u(u)=E\phi(E)$, and $\sigma_u(u)=\sigma(E(u))$.
For one elastic collision, $\Delta u=u'-u$:
\[
 p(\Delta u)=\frac{e^{-\Delta u}}{1-\alpha},
 \quad 0\le\Delta u\le\ln(1/\alpha),\qquad\EE[\Delta u]=\xi.
\]
Bin flux is \emph{integrated}: plot $\phi_g/\Delta E_g$ per eV or
$\phi_g/\Delta u_g$ per unit lethargy.
\RefLabel{Thermal equilibrium.} Set $\theta=k_BT$. On $E>0$,
\[
 p_n(E)=\frac{2\sqrt E}{\sqrt\pi\theta^{3/2}}e^{-E/\theta},
 \qquad w_T(E)=\frac{E}{\theta^2}e^{-E/\theta}.
\]
$p_n$: normalized \emph{neutron density}; $w_T$: normalized \emph{flux}.
Room temperature: $\theta\simeq0.0253$ eV.
\[
 \bar v=\sqrt{\frac{8\theta}{\pi m_n}},\quad
 \langle E\rangle_n=\tfrac32\theta,\quad
 \tfrac12m_n\bar v^2=\frac4\pi\theta\ne\langle E\rangle_n.
\]
Useful integrals: $\int_0^\infty Ee^{-E/\theta}\dd E=\theta^2$;
$\int_0^\infty\sqrt E e^{-E/\theta}\dd E=\sqrt\pi\theta^{3/2}/2$.

\RefHead{7}{Fast, slowing-down, and resonance spectra}{L6,9--10}
\RefLabel{Fission births.} $\int\chi\dd E=1$; $\chi$ need not equal
$\phi/\phi_{\rm tot}$. Course Watt model ($E$ in \textbf{MeV}):
\[
 \chi(E)=C_\chi e^{-1.036E}\sinh\sqrt{2.29E},\qquad
 C_\chi\simeq0.452463\unit{MeV^{-1}}.
\]
Ignoring in-scatter gives $\phi_F\propto\chi/\Sigma_t$;
only a constant $\Sigma_t$ gives $\phi_F\propto\chi$.
\RefLabel{Slowing down below the source.} No absorption or local births;
$E'$ is now \emph{incident}, $E$ \emph{outgoing}:
\[
 \Sigma_s(E)\phi(E)=\int_E^{E/\alpha}
 \frac{\Sigma_s(E')\phi(E')}{(1-\alpha)E'}\dd E',
\]
\[
 \Sigma_s(E)\phi(E)=\frac{C}{E},\quad q=\xi C,
 \quad \phi(E)=\frac{q}{\xi\Sigma_s(E)E}.
\]
For H-1, the upper limit is $\infty$. $q$ is the slowing-down density
(cm$^{-3}$s$^{-1}$). A $1/E$ \emph{flux} requires constant $\Sigma_s$.
For mixtures with constant scattering, use $\sum_i\xi_i\Sigma_s^i$.

\vspace{0.25cm}

\RefLabel{Resonant fuel $f$, moderator $m$.} Background $C/E$;
$\Sigma_{s,0}=\Sigma_{s,0}^m+\Sigma_{s,0}^f$ is background scattering.
\[
 \phi_{\rm NR}\simeq\frac{C}{\Sigma_t(E) E},\qquad
 \phi_{\rm WR}\simeq\frac{C}{(\Sigma_t(E)-\Sigma_s^f)E}.
\]
NR: $\Delta E_{\rm res}\ll(1-\alpha_i)E_r$ for both species.
WR: $A_f\to\infty$ cancels same-energy fuel scattering; the moderator
remains NR. Resonances depress $\phi$ (\textbf{self-shielding}), but the
absorption rate is $\Sigma_a\phi$. Use a common normalization.

\RefHead{8}{Useful course data}{L1,8--9,12}

\RefLabel{Atomic masses (u); retain digits until the final $Q$.}
\par\vspace{2pt}{\centering\fontsize{8.0}{9.2}\selectfont
\setlength{\tabcolsep}{2pt}
\begin{tabular}{@{}lr@{\quad}lr@{\quad}lr@{}}
\is{1}{H}&1.007825031898&\is{4}{He}&4.00260325413&\is{7}{Li}&7.01600343426\\
\is{9}{Be}&9.012183062&\is{10}{B}&10.012936862&\is{12}{B}&12.014352638\\
\is{13}{C}&13.00335483534&\is{60}{Co}&59.933815536&\is{60}{Ni}&59.930785129
\end{tabular}\par}\vspace{2pt}

\par\vspace{2pt}{\centering\fontsize{8.0}{9.2}\selectfont
\begin{tabular}{@{}lrrr@{}}
\toprule
Moderator & $\xi$ & $\xi\Sigma_s$ (cm$^{-1}$) & $\xi\Sigma_s/\Sigma_a$\\
\midrule
H$_2$O&0.93&1.28&58\\
D$_2$O&0.51&0.18&21000\\
C&0.158&0.056&200\\
\bottomrule
\end{tabular}
\par}\vspace{2pt}
FNRP Table 3.1: moderator ratio uses \emph{thermal} absorption.

\par\vspace{2pt}
{\centering\fontsize{8}{9.2}\selectfont
\begin{tabular}{@{}lrrrrrr@{}}
\toprule
 & $\nu_T$ & $\sigma_{f,T}$ & $\sigma_{a,T}$ & $I_a$ & $\sigma_{f,F}$ & $\sigma_{a,F}$\\
\midrule
U-235 & 2.43 & 505 & 591 & 404 & 1.22 & 1.29\\
U-238 & --- & $\simeq0$ & 2.42 & 278 & 0.304 & 0.361\\
Pu-239 & 2.90 & 698 & 973 & 474 & 1.81 & 1.86\\
\bottomrule
\end{tabular}\par}
$T$: room-temperature Maxwellian \emph{flux averages}, not 0.0253-eV
point values; $F$: fission-spectrum averages; $I_a$: 1--$10^5$ eV.

\switchcolumn\RaggedRight
% ---------- PAGE 2, RIGHT ----------
\RefHead{9}{Rate-preserving energy averages}{L8,12--13}
Effective cross sections:
\[
 \bar\Sigma_x=\frac{\int\Sigma_x\phi\dd E}{\int\phi\dd E},\quad
 \overline{\nu\Sigma_f}=\frac{\int\nu\Sigma_f\phi\dd E}{\int\phi\dd E}.
\]
Average $\nu\Sigma_f$ as a product, not as independent averages.
\[
 I_x=\int_{E_l}^{E_u}\frac{\sigma_x(E)}{E}\dd E,\qquad
 \bar\sigma_{x,I}=\frac{I_x}{\ln(E_u/E_l)}\quad(1/E\text{ weight}).
\]
NR weighting instead uses $1/[\Sigma_t(E)E]$ in \emph{both} integrals.
For $1/v$ cross sections and Maxwellian flux,
\[
 \bar\sigma_{x,T}=\frac{\sqrt\pi}{2}\sigma_x(E_*)
 \sqrt{\frac{E_*}{k_BT}}.
\]


\RefHead{10}{Groups, transfer, and condensation}{L13}
Number from high to low energy: $E_0>\cdots>E_G$,
$\phi_g=\int_{E_g}^{E_{g-1}}\phi(E)\dd E$.
Course bounds (eV): $(10^7,10^5,1,10^{-3})$ for $(F,I,T)$;
$\Delta u_g=\ln(E_{g-1}/E_g)$.
Apply the flux average above within each group:
$\Sigma_{x,g}=\int_g\Sigma_x\phi\dd E/\phi_g$;
$\chi_g=\int_g\chi\dd E$, $s_g=\int_gs\dd E$.

\vspace{0.25cm}

\RefLabel{Destination $g\gets$ incident $g'$.} For elastic scattering,
\[
 P_g^i(E')=\frac{[\min(E_{g-1},E')-\max(E_g,\alpha_i E')]_+}
                  {(1-\alpha_i)E'},\quad[z]_+=\max(z,0),
\]
\[
 \Sigma_{s,g\gets g'}=\sum_i n_i
 \frac{\int_{g'}\sigma_s^i(E')P_g^i(E')\phi(E')\dd E'}{\phi_{g'}}.
\]
The denominator is the \emph{incident}-group flux ($\propto\Delta u_{g'}$
for $1/E'$ weighting). No upscatter: entries above the diagonal vanish.

\RefLabel{Collapse fine groups into a coarse group $G$.}
\[
 \phi_G=\sum_{g\in G}\phi_g,\qquad
 \Sigma_{x,G}=\frac{\sum_{g\in G}\Sigma_{x,g}\phi_g}{\phi_G}.
\]
Add rates, not cross sections!



\RefHead{11}{Infinite-medium neutron balance}{L12--14}
Steady, homogeneous, no leakage; number-conserving scattering.
For each energy group,
\[
 \Sigma_{t,g}\phi_g
 =
 \sum_{g'}\Sigma_{s,g\gets g'}\phi_{g'}
 +\chi_g\sum_{g'}(\nu\Sigma_f)_{g'}\phi_{g'}
 +s_g .
\]
Thus, scattering out of group $g$ is included in $\Sigma_{t,g}\phi_g$,
while scattering into $g$ appears on the right.
For a fixed source,
\[
 \sum_{g'}
 \left[
   \Sigma_{t,g}\delta_{gg'}
   -\Sigma_{s,g\gets g'}
   -\chi_g(\nu\Sigma_f)_{g'}
 \right]\phi_{g'}
 =s_g .
\]
For the eigenvalue problem,
\[
 \sum_{g'}
 \left[
   \Sigma_{t,g}\delta_{gg'}
   -\Sigma_{s,g\gets g'}
 \right]\phi_{g'}
 =
 \frac{\chi_g}{k_{\text{eff}}}
 \sum_{g'}(\nu\Sigma_f)_{g'}\phi_{g'} .
\]
Only fission is divided by $k_{\text{eff}}$. The dominant eigenvector
sets shape, not power; e.g., $\sum_g\widehat{\phi}_g=1$.
A physical multiplying fixed-source solution requires $\nu\Sigma_f/\Sigma_a \equiv k_\infty<1$!!

\RefLabel{One group.}
\[
  \Sigma_a \phi= \frac{1}{k_{\text{eff}}} \nu\Sigma_f  \phi + s
  \Longrightarrow
  \phi = \frac{s}{\Sigma_a-\nu\Sigma_f/k_{\text{eff}}}
  = \frac{s/\Sigma_a}{1-k_{\infty}/k_{\text{eff}}}
\]

\RefLabel{Two groups, no upscatter.}
Let $(\chi_1,\chi_2)=(1,0)$, $(s_1,s_2)=(s_1,0)$. Then
\[
 \Sigma_{a,2}\phi_2
 =
 \Sigma_{s,2\gets1}\phi_1,
 \qquad
 \frac{\phi_2}{\phi_1}
 =
 \frac{\Sigma_{s,2\gets1}}{\Sigma_{a,2}},
\]
and
\[
 k_\infty
 =
 \frac{
   (\nu\Sigma_f)_1\phi_1+(\nu\Sigma_f)_2\phi_2
 }{
   \left(\Sigma_{a,1}+\Sigma_{s,2\gets1}\right)\phi_1
 }
 =
 \frac{
   (\nu\Sigma_f)_1
   +(\nu\Sigma_f)_2(\phi_2/\phi_1)
 }{
   \Sigma_{a,1}+\Sigma_{s,2\gets1}
 }.
\]
For a fixed fast-group source,
\[
 \phi_1
 =
 \frac{s_1}
 {\left(\Sigma_{a,1}+\Sigma_{s,2\gets1}\right)(1-k_\infty)},
 \qquad
 \phi_2
 =
 \frac{\Sigma_{s,2\gets1}}{\Sigma_{a,2}}\phi_1 .
\]


\end{handoutcolumns}
\end{document}
